% Part: first-order-logic
% Chapter: sequent-calculus
% Section: proof-theoretic-notions

\documentclass[../../../include/open-logic-section]{subfiles}

\begin{document}

\begin{editorial}
  या विभागात नैसर्गिक निगमनासाठी सिद्धतायोग्यता संबंध आणि सुसंगतता यांच्या
  व्याख्या एकत्रित केल्या आहेत.
\end{editorial}

\iftag{FOL}
      {\olfileid{fol}{seq}{ptn}}
      {\olfileid{pl}{seq}{ptn}}

\olsection{सिद्धता-उपपत्तीय संकल्पना}

\begin{explain}
जशा आपण अनेक महत्त्वाच्या चिन्हार्थविषयक संकल्पना (वैधता, तार्किक निष्पन्नता,
पूर्ततायोग्यता) परिभाषित केल्या आहेत, तशाच त्यांना अनुरूप
\emph{सिद्धता-उपपत्तीय संकल्पना} आता परिभाषित करू. !!{structure}s मध्ये
!!{sentence}s ची पूर्ति होते की नाही, याचा आधार घेऊन या संकल्पना परिभाषित
केल्या जात नाहीत; त्याऐवजी ठरावीक क्रमवर्तींची !!{derivability} किंवा
!!{nonderivability} यांचा आधार घेतला जातो. या दोन प्रकारच्या संकल्पना
एकमेकांशी जुळतात, हा एक महत्त्वाचा शोध होता. त्या जुळतात, हाच
\emph{निर्दोषता प्रमेय} आणि \emph{संपूर्णता प्रमेय} यांचा आशय आहे.
\end{explain}

\begin{defn}[प्रमेये]
!!^a{sentence}~$!A$ हे \emph{प्रमेय} असते, जर $\Log{LK}$ मध्ये
$\quad \Sequent !A$ या क्रमवर्तीची !!a{derivation} असेल. $!A$ हे प्रमेय
असेल तर आपण $\Proves !A$ लिहितो आणि ते प्रमेय नसेल तर $\Proves/ !A$ लिहितो.
\end{defn}

\begin{defn}[!!^{derivability}]
!!^a{sentence}~$!A$ हे !!{sentence}s च्या संच~$\Gamma$ \emph{पासून
!!{derivable}} असते, म्हणजे $\Gamma \Proves !A$, तेव्हा आणि केवळ तेव्हाच,
जेव्हा एखादा सांत उपसंच~$\Gamma_0 \subseteq \Gamma$ आणि $\Gamma_0$ मधील
!!{sentence}s ची एखादी क्रमिका $\Gamma_0'$ अशी असते की $\Log{LK}$
$\Gamma_0' \Sequent !A$ हे !!{derive}s करते. $!A$ हे $\Gamma$ पासून
!!{derivable} नसेल, तर आपण $\Gamma \Proves/ !A$ लिहितो.
\end{defn}

आकुंचन, दुर्बलीकरण आणि अदलाबदल नियमांमुळे $\Gamma_0'$ मधील
!!{sentence}s चा क्रम आणि त्यांची पुनरावृत्ती कितीदा होते हे महत्त्वाचे नसते:
$\Gamma_0' \Sequent !A$ ही क्रमवर्ती !!{derivable} असेल, तर
$\Gamma_0'$ मध्ये असलेली तीच !!{sentence}s असणाऱ्या कोणत्याही
$\Gamma_0''$ साठी $\Gamma_0'' \Sequent !A$ ही क्रमवर्तीसुद्धा निष्पन्न करता
येते. उदाहरणार्थ, $\Gamma_0 = \{!B, !C\}$ असेल, तर
$\Gamma_0' = \tuple{!B, !B, !C}$ आणि $\Gamma_0'' =
\tuple{!C, !C, !B}$ या दोन्ही क्रमिकांमध्ये $\Gamma_0$ मधीलच
!!{sentence}s आहेत. यांपैकी एक क्रमिका असणारी क्रमवर्ती !!{derivable} असेल,
तर दुसरी असणारी क्रमवर्तीसुद्धा निष्पन्न करता येते; उदा.:
\begin{prooftree}
  \AxiomC{}
  \Deduce$!B, !B, !C \fCenter !A$
  \RightLabel{\LeftR{\Contraction}}
  \UnaryInf$!B, !C \fCenter !A$
  \RightLabel{\LeftR{\Exchange}}
  \UnaryInf$!C, !B \fCenter !A$
  \RightLabel{\LeftR{\Weakening}}
  \UnaryInf$!C, !C, !B \fCenter !A$
\end{prooftree}
यापुढे $\Gamma_0$ हा !!{sentence}s चा सांत संच असेल, तर
$\Gamma_0 \Sequent !A$ याचा अर्थ असा कोणताही क्रमवर्ती घेऊ की ज्याच्या
पूर्वांगात $\Gamma_0$ मधील !!{sentence}s ची क्रमिका आहे; तसेच आवश्यक असल्यास
आकुंचन, अदलाबदल आणि दुर्बलीकरणे गृहीत धरू.

\begin{defn}[सुसंगतता]
वाक्यांचा संच~$\Gamma$ हा \emph{विसंगत} असतो तेव्हा आणि केवळ तेव्हाच, जेव्हा
एखादा सांत उपसंच~$\Gamma_0 \subseteq \Gamma$ असा असतो की $\Log{LK}$
$\Gamma_0 \Sequent \quad$ हे !!{derive}s करते. $\Gamma$ विसंगत नसेल,
म्हणजे प्रत्येक सांत $\Gamma_0 \subseteq \Gamma$ साठी $\Log{LK}$
$\Gamma_0 \Sequent \quad$ हे !!{derive} करत नसेल, तर त्याला
\emph{सुसंगत} म्हणतो.
\end{defn}

\begin{prop}[परावर्तिता]
\ollabel{prop:reflexivity}
$!A \in \Gamma$ असेल, तर $\Gamma \Proves !A$.
\end{prop}

\begin{proof}
$!A \Sequent !A$ ही आरंभीची क्रमवर्ती !!{derivable} आहे आणि $\{!A\}
\subseteq \Gamma$.
\end{proof}
  
\begin{prop}[एकस्वनिकता]
\ollabel{prop:monotonicity}
$\Gamma \subseteq \Delta$ आणि $\Gamma \Proves !A$ असेल, तर $\Delta
\Proves !A$.
\end{prop}

\begin{proof}
$\Gamma \Proves !A$ असे समजा; म्हणजे, एखादा सांत $\Gamma_0
\subseteq \Gamma$ असा आहे की $\Gamma_0 \Sequent !A$ ही क्रमवर्ती
!!{derivable} आहे. $\Gamma \subseteq \Delta$ असल्यामुळे $\Gamma_0$ हा
$\Delta$ चासुद्धा सांत उपसंच आहे. म्हणून $\Gamma_0 \Sequent !A$ ची
!!{derivation} ही $\Delta \Proves !A$ हेसुद्धा दाखवते.
\end{proof}

\begin{prop}[संक्रमकता]
\ollabel{prop:transitivity}
$\Gamma \Proves !A$ आणि $\{!A\} \cup \Delta \Proves
!B$ असेल, तर $\Gamma \cup \Delta \Proves !B$.
\end{prop}

\begin{proof}
$\Gamma \Proves !A$ असेल, तर एखादा सांत $\Gamma_0 \subseteq \Gamma$
आणि $\Gamma_0 \Sequent !A$ ची !!a{derivation}~$\pi_0$ असते.
$\{!A\} \cup \Delta \Proves !B$ असेल, तर एखाद्या सांत उपसंचासाठी
$\Delta_0 \subseteq \Delta$, $!A, \Delta_0 \Sequent !B$ ची
!!a{derivation}~$\pi_1$ असते. पुढील !!{derivation} विचारात घ्या:
\begin{prooftree}
  \AxiomC{}
  \RightLabel{$\pi_0$}
  \Deduce$\Gamma_0 \fCenter !A$
  \AxiomC{}
  \RightLabel{$\pi_1$}
  \Deduce$!A, \Delta_0 \fCenter !B$
  \RightLabel{\Cut}
  \BinaryInf$\Gamma_0, \Delta_0 \fCenter !B$
\end{prooftree}
$\Gamma_0 \cup \Delta_0 \subseteq \Gamma \cup \Delta$ असल्यामुळे यावरून
$\Gamma \cup \Delta \Proves !B$ हे दिसते.
\end{proof}

विशेषतः, याचा अर्थ असा की $\Gamma \Proves !A$ आणि $!A
\Proves !B$ असेल, तर $\Gamma \Proves !B$. तसेच $!A_1,
\dots, !A_n \Proves !B$ असेल आणि प्रत्येक~$i$ साठी $\Gamma \Proves !A_i$
असेल, तर $\Gamma \Proves !B$, हेसुद्धा यावरून मिळते.

\begin{prop}
\ollabel{prop:incons}
$\Gamma$ विसंगत असतो तेव्हा आणि केवळ तेव्हाच, जेव्हा प्रत्येक
  वाक्य~$!A$ साठी $\Gamma \Proves {!A}$.
\end{prop}

\begin{proof}
सराव.
\end{proof}

\begin{prob}
\olref[fol][seq][ptn]{prop:incons} सिद्ध करा.
\end{prob}

\begin{prop}[संहतता]
  \ollabel{prop:proves-compact}
  \begin{enumerate}
  \item $\Gamma \Proves !A$ असेल, तर एखादा सांत उपसंच $\Gamma_0
    \subseteq \Gamma$ असा आहे की $\Gamma_0 \Proves !A$.
  \item $\Gamma$ चा प्रत्येक सांत उपसंच सुसंगत असेल, तर $\Gamma$ सुसंगत
    असतो.
  \end{enumerate}
\end{prop}

\begin{proof}
  \begin{enumerate}
    \item $\Gamma \Proves !A$ असेल, तर एखादा सांत उपसंच
      $\Gamma_0 \subseteq \Gamma$ असा आहे की $\Gamma_0
      \Sequent !A$ या क्रमवर्तीची !!a{derivation} आहे. परिणामी,
      $\Gamma_0 \Proves !A$.
    \item $\Gamma$ विसंगत असेल, तर एखादा सांत
      उपसंच~$\Gamma_0 \subseteq \Gamma$ असा आहे की $\Log{LK}$
      $\Gamma_0 \Sequent \quad$ हे !!{derive}s करते. पण मग $\Gamma_0$ हा
      $\Gamma$ चा विसंगत सांत उपसंच आहे.
  \end{enumerate}
\end{proof}

\end{document}
